% ============================================================
%  equations_expo.tex
%  Ecuaciones principales del fuzzy RDD — listo para Overleaf
%  Pensión 65 × Billetera Digital, ENAHO 2024
% ============================================================

% ── Notación ──────────────────────────────────────────────
% i       = individuo
% D_i     = 1 si la persona i recibe Pensión 65 (P5567A = 1)
% Y_i     = TIENE_BILLETERA (billetera digital, tenencia o uso)
% Z_i     = 1(Edad_i >= 65)   [instrumento / discontinuidad]
% r_i     = Edad_i - 65       [running variable centrada]
% h       = ancho de banda (h* = 14 años)
% K(·)    = kernel triangular
% ──────────────────────────────────────────────────────────

% ── Ecuación 1: Primera Etapa (First Stage) ───────────────
% Estima cómo el instrumento Z_i predice el tratamiento D_i.
% El coeficiente rho (ρ) es el salto en la probabilidad de
% recibir Pensión 65 exactamente en el corte de edad 65.

\begin{equation}
    \label{eq:first_stage}
    D_i
    \;=\;
    \alpha_0
    + \underbrace{\rho}_{\text{1ra etapa}}
        \cdot \mathbf{1}[\text{Edad}_i \geq 65]
    + \alpha_1 \, r_i
    + \alpha_2 \, r_i \cdot \mathbf{1}[r_i \geq 0]
    + \varepsilon_i
    \quad
    \left| \, |r_i| \leq h \right.
\end{equation}

% Resultado estimado (Candidato A, muestra completa, h = 14):
%   \hat{\rho} = 0.104,  F-estadístico \approx 310  (instrumento fuerte)

% ── Ecuación 2: Segunda Etapa (IV / 2SLS) ────────────────
% Regresa el outcome Y_i sobre la predicción \hat{D}_i
% de la primera etapa. El coeficiente tau (τ) es el LATE:
% Local Average Treatment Effect para los compliers.

\begin{equation}
    \label{eq:second_stage}
    Y_i
    \;=\;
    \beta_0
    + \underbrace{\tau}_{\text{LATE}}
        \cdot \widehat{D}_i
    + \beta_1 \, r_i
    + \beta_2 \, r_i \cdot \mathbf{1}[r_i \geq 0]
    + u_i
    \quad
    \left| \, |r_i| \leq h \right.
\end{equation}

% Resultado estimado (Candidato A, muestra completa):
%   \hat{\tau}_{LATE} = -0.012,  IC_{95\%} = [-0.190,\; +0.166]
%   No rechazamos nulo. Cota superior excluye efectos > +17 pp.

% ── Ecuación 3: Relación ITT – LATE (Wald / 2SLS) ────────
% La estimación de variables instrumentales en el RDD se
% interpreta como el cociente entre la forma reducida (RF,
% efecto del instrumento sobre el outcome) y la primera
% etapa (FS, efecto del instrumento sobre el tratamiento).
% Siguiendo Angrist y Pischke (2009), cap. 4.

\begin{equation}
    \label{eq:late_wald}
    \hat{\tau}_{\text{LATE}}
    \;=\;
    \frac{\hat{\tau}_{\text{RF}}}{\hat{\rho}}
    \;=\;
    \frac{
        \displaystyle\lim_{r \downarrow 0} \mathbb{E}[Y_i \mid r_i = r]
        -
        \displaystyle\lim_{r \uparrow 0}  \mathbb{E}[Y_i \mid r_i = r]
    }{
        \displaystyle\lim_{r \downarrow 0} \mathbb{E}[D_i \mid r_i = r]
        -
        \displaystyle\lim_{r \uparrow 0}  \mathbb{E}[D_i \mid r_i = r]
    }
\end{equation}

% El LATE identifica el efecto causal de Pensión 65 sobre
% la adopción de billetera digital ÚNICAMENTE para los
% "compliers": individuos que reciben el bono si y solo si
% tienen 65 años o más (Imbens y Angrist, 1994).
%
% Referencia:
%   Angrist, J. D., & Pischke, J. S. (2009).
%   \textit{Mostly Harmless Econometrics}.
%   Princeton University Press. Cap. 4.6.
